\documentclass{amsart}
\usepackage{amsmath, amssymb}
\title{Math 21b --- Lesson Plan: Week 1}
\subtitle{Introduction to Linear Systems}
\date{}
\begin{document}
\maketitle

\textbf{Learning Objectives.} By the end of this week, students will be able to:
\begin{itemize}
    \item Set up and interpret a system of linear equations.
    \item Use row operations to reduce an augmented matrix to echelon form.
    \item Identify pivot positions, free variables, and the solution set of a system.
\end{itemize}

\bigskip

\textbf{Lecture 1 (50 min) --- Systems of Linear Equations}

\textit{Minutes 0--10: Motivation.} Open with a concrete example: fitting a line through data points, or balancing a chemical equation. Ask: ``What does it mean to solve a system?'' Show that three planes in $\mathbb{R}^3$ can intersect in a point, a line, or not at all.

\textit{Minutes 10--30: Row Operations and Echelon Form.} Define elementary row operations. Work through a $3\times 3$ system step by step on the board; have students follow on paper. Define echelon form and RREF.

\textit{Minutes 30--45: Back Substitution and Parametric Form.} Demonstrate back substitution. Introduce parametric vector form for systems with free variables.

\textit{Minutes 45--50: Preview.} ``Next time: we'll reinterpret all of this using matrices and linear transformations.''

\bigskip

\textbf{Lecture 2 (50 min) --- Matrix Notation and Gaussian Elimination}

\textit{Minutes 0--5: Review quiz.} Two quick true/false questions on row operations.

\textit{Minutes 5--25: Matrix Notation.} Define coefficient matrix and augmented matrix. Introduce matrix--vector product $Ax$. Rewrite $Ax = b$.

\textit{Minutes 25--40: Gaussian Elimination Algorithm.} State the algorithm formally; work a $4\times 4$ example. Discuss pivot columns vs.\ free columns.

\textit{Minutes 40--50: Group Activity.} Students work in pairs on a system with a parameter $k$ and determine for which $k$ the system has 0, 1, or $\infty$ solutions.

\bigskip

\textbf{Section (50 min)}
\begin{itemize}
    \item Worksheet 1 (Problems 1--2): row reduction practice.
    \item Discussion: what does it mean geometrically to have infinitely many solutions?
    \item Preview of Problem Set 1.
\end{itemize}

\bigskip

\textbf{Assigned Reading.} Lay \S1.1--1.3 (or equivalent).

\textbf{Assigned Work.} Problem Set 1, Problems 1 and 4 (due next week).

\end{document}
